\begin{mvtable}
\centering
{\small 续表}\\[2pt]
\begin{tabular}{|c|c||c|c|}
\hline
$x_{1}$(硬度) & $x_{2}$(弯曲强度) & $x_{1}$(硬度) & $x_{2}$(弯曲强度)\\
\hline
$1\,752$ & $9\,469$ & $1\,587$ & $8\,309$\\
$2\,067$ & $6\,410$ & $2\,208$ & $9\,559$\\
$2\,365$ & $10\,327$ & $1\,487$ & $6\,255$\\
$1\,646$ & $7\,320$ & $2\,206$ & $10\,723$\\
$1\,579$ & $8\,196$ & $2\,332$ & $5\,430$\\
$1\,880$ & $9\,709$ & $2\,540$ & $12\,090$\\
$1\,773$ & $10\,370$ & $2\,322$ & $10\,072$\\
\hline
\end{tabular}\\[3pt]
{\footnotesize 资料来源: 数据经美国林业制品实验室许可使用.}
\end{mvtable}

\noindent (a) 构造并画出 $[\mu_{1},\mu_{2}]'$ 的 $95\%$ 置信椭圆, 其中 $\mu_{1}=E(X_{1}),\mu_{2}=E(X_{2})$.

\noindent (b) 假设硬度及弯曲强度的“标准”值分别为 $\mu_{10}=2\,000$ 及 $\mu_{20}=10\,000$, 已知 (a) 的结论, 讨论\hyperref[tab:5.11]{表 5.11} 中的数据是否与这些数值相容.

\noindent (c) 用 Q-Q 图及散布图来解释二维正态分布是否为该总体的一个可行的模型.

\medskip
\noindent\mvbold{5.20}\quad 一位野生动物生态学家对一个由 $n=45$ 只雄性钩镰鸢组成的样本测量了两项指标: $x_{1}=$ 尾长 (毫米), $x_{2}=$ 翅长 (毫米), 数据列于\hyperref[tab:5.12]{表 5.12}.

\begin{mvtable}
\phantomsection\label{tab:5.12}\mvtabcap{表 5.12\quad 钩镰鸢数据}
\centering
{\small
\begin{tabular}{|c|c|c|c|c|c|}
\hline
$x_{1}$ & $x_{2}$ & $x_{1}$ & $x_{2}$ & $x_{1}$ & $x_{2}$\\
(尾长) & (翅长) & (尾长) & (翅长) & (尾长) & (翅长)\\
\hline
$191$ & $284$ & $186$ & $266$ & $173$ & $271$\\
$197$ & $285$ & $197$ & $285$ & $194$ & $280$\\
$208$ & $288$ & $201$ & $295$ & $198$ & $300$\\
$180$ & $273$ & $190$ & $282$ & $180$ & $272$\\
$180$ & $275$ & $209$ & $305$ & $190$ & $292$\\
$188$ & $280$ & $187$ & $285$ & $191$ & $286$\\
$210$ & $283$ & $207$ & $297$ & $196$ & $285$\\
$196$ & $288$ & $178$ & $268$ & $207$ & $286$\\
$191$ & $271$ & $202$ & $271$ & $209$ & $303$\\
$179$ & $257$ & $205$ & $285$ & $179$ & $261$\\
$208$ & $289$ & $190$ & $280$ & $186$ & $262$\\
$202$ & $285$ & $189$ & $277$ & $174$ & $245$\\
$200$ & $272$ & $211$ & $310$ & $181$ & $250$\\
$192$ & $282$ & $216$ & $305$ & $189$ & $262$\\
$199$ & $280$ & $189$ & $274$ & $188$ & $258$\\
\hline
\end{tabular}}\\[3pt]
{\footnotesize 资料来源: 数据经 S. Temple 许可使用.}
\end{mvtable}

利用表中的数据:

\noindent (a) 构造并画出总体均值 $\mu_{1}$ 和 $\mu_{2}$ 的 $95\%$ 置信椭圆, 假定雄性钩镰鸢的 $\mu_{1}=190$ 毫米, $\mu_{2}=275$ 毫米. 试问这些值是否为雄性钩镰鸢的平均尾长和平均翅膀长的真值? 请予以解释.
